\chapter*{Start here: from zero}
\addcontentsline{toc}{chapter}{Start here: from zero}
\markboth{Start here}{Start here}

Machine learning means getting a computer to find a rule from examples instead
of being told the rule. You do not need to have used it before. You do need two
small things: the ability to read a short Python program, and a feel for vectors
and matrices. This chapter gives both in a few pages, using programs you can run.

\section*{Python in one page}

\lstinputlisting[caption={\texttt{code/start/python\_primer.py}: values, a list, a function, a loop.},label={lst:start-py}]
{code/start/python_primer.py}

Running it prints:
\begin{lstlisting}[numbers=none]
count   = 4
average = 116.25
largest = 140
squares = [14400, 9025]
\end{lstlisting}
(Check the average by hand: $120+95+140+110=465$, and $465/4=116.25$.)
A \emph{list} is an ordered row of values; \texttt{def} defines a function;
indentation, not braces, marks which lines belong to the loop or function; and
\texttt{prices[:2]} means ``the first two elements.'' That is enough Python to
start Chapter~1.

\section*{Vectors and matrices in one page}

A \emph{vector} is a list of numbers, such as the measurements of one penguin:
$\mathbf x=(39.1,\ 18.7,\ 181)$. A \emph{matrix} is a table of numbers: stack one
vector per row and you have a dataset $X$ with $n$ rows (samples) and $d$ columns
(features). Multiplying a matrix by a vector of weights $\mathbf w$ gives one
number per row, the prediction for each sample; this single operation is the
heart of a linear model. In NumPy:

\lstinputlisting[caption={\texttt{code/start/numpy\_primer.py}: a dataset as a matrix, and a prediction for each row.},label={lst:start-np}]
{code/start/numpy_primer.py}

\begin{lstlisting}[numbers=none]
shape of X: (3, 2)
X @ w     : [12. 34. 56.]
column means: [3. 4.]
\end{lstlisting}
Check the first prediction by hand: $1\cdot10+2\cdot1=12$.

This is all the linear algebra you need to \emph{start}. When the book needs more
(least squares, SVD, gradients), the \emph{Before you start} box at the top of the
chapter points to the exact chapter of \textbf{Book~3} and restates the idea in a
sentence, so you can read on without it.

\section*{How the books of this series fit together}

\begin{center}
\begin{tabular}{P{2.4cm}P{9.0cm}}
\toprule
\thead{Book} & \thead{Connection}\\
\midrule
3 Linear algebra & The mathematical spine of this book. We cite it by chapter: least squares (Ch.~8), PCA/SVD (Ch.~10), gradients (Ch.~11), optimisation (Ch.~14), entropy (Ch.~15).\\
1 Competitive programming & Where to learn algorithmic thinking and cost estimation; helpful for Chapters~14 and~18 here.\\
6 Quant research & Applies the evaluation discipline of Chapter~3 here to financial data.\\
\bottomrule
\end{tabular}
\end{center}

\section*{How to use the boxes in each chapter}

Each chapter opens with \emph{After this chapter you can}, naming what you can
\emph{do}. From Chapter~2, a \emph{Before you start} box lists what the chapter
relies on. Do the ``Drill'' and ``Exercises'' sections: reading a training curve
is a skill, and skills need practice.
